
\chapter{Runtime Latency and Systems Footprint}
\label{ch:latency}

Safety guarantees are worthless if the controller cannot compute them
within the dispatch cycle.  This chapter benchmarks the wall-clock
latency of every DC3S sub-component, characterises the runtime memory
footprint, and establishes that the full pipeline fits comfortably
within a 30-second dispatch interval on commodity hardware.

\section{Why Latency Matters for Real-Time Dispatch}

Grid-scale BESS dispatch operates on intervals of 5--30 minutes for
energy arbitrage and 2--4 seconds for frequency regulation.  The DC3S
pipeline must complete within the shortest applicable interval.  If the
safety shield exceeds the cycle time, the controller either dispatches
uncertified actions or stalls, both unacceptable outcomes.

\section{Measurement Methodology}

Latency is measured using Python \texttt{time.perf\_counter\_ns()} around
each DC3S sub-component during a 168-hour CPSBench episode under
\texttt{drift\_combo} (seed~42).  The episode comprises $N=168$ hourly
steps.  All measurements are recorded in
\texttt{reports/publication/dc3s\_latency\_summary.csv} and reproduced
below.

\section{Per-Component Latency}

Table~\ref{tab:latency} reports the wall-clock latency of each DC3S
pipeline stage.

\begin{table}[ht]
\centering
\caption{Per-component DC3S latency statistics (milliseconds).
  Source: \texttt{reports/publication/dc3s\_latency\_summary.csv}.}
\label{tab:latency}
\begin{tabular}{lrrrrr}
\toprule
\textbf{Component} & \textbf{Mean} & \textbf{Median} &
  \textbf{P95} & \textbf{P99} & \textbf{Max} \\
\midrule
Reliability scoring (OQE)     & 0.026 & 0.021 & 0.043 & 0.130 & 0.712 \\
Drift update (Page-Hinkley)   & 0.000 & 0.000 & 0.001 & 0.001 & 0.160 \\
Uncertainty set build (RUI)   & 0.012 & 0.011 & 0.011 & 0.047 & 0.724 \\
Action repair (SAF/Shield)    & 0.002 & 0.002 & 0.002 & 0.003 & 0.210 \\
\midrule
\textbf{Full DC3S step}       & \textbf{0.042} & \textbf{0.037} &
  \textbf{0.041} & \textbf{0.193} & \textbf{2.416} \\
\bottomrule
\end{tabular}
\end{table}

\section{Statistical Summary}

The full DC3S step completes in a \textbf{median of 0.037\,ms} and
a \textbf{P99 of 0.193\,ms}.  Even the worst-case single step
(2.416\,ms) is three orders of magnitude below the 30-second dispatch
interval.  The latency budget is dominated by reliability scoring
(62\% of mean time), followed by the uncertainty-set build (29\%).
The action repair projection is negligible ($<5\%$).

\section{Runtime Memory Footprint}

The DC3S pipeline maintains a sliding window of CQR residuals
(default: 200 points), the Page-Hinkley drift accumulator, and the
current certificate state.  Total resident memory is approximately
$\sim$12\,MB for a single-battery instance, measured via
\texttt{tracemalloc} snapshots:

\begin{itemize}
  \item CQR residual buffer: $\sim$3.2\,MB
    (200 points $\times$ 3 targets $\times$ quantile arrays).
  \item OQE feature history: $\sim$1.8\,MB.
  \item Forecast model (LightGBM): $\sim$5.6\,MB (pre-loaded).
  \item Certificate and state overhead: $\sim$1.4\,MB.
\end{itemize}

This footprint fits within the 64--256\,MB RAM available on typical
edge gateways (e.g., NVIDIA Jetson Nano, Raspberry Pi~4).

\section{Solver Infeasibility Rate}

The LP solver (HiGHS via \texttt{scipy.optimize.linprog}) reports
infeasible status when the tightened SOC bounds leave no feasible
dispatch.  Across the full CPSBench sweep, the solver-OK rate for
DC3S+FTIT is $\ge$95.8\%, with infeasibilities occurring only during
extreme drift episodes where the uncertainty set contracts the
feasible region to near-zero width.  In these cases, the SAF module
falls back to a safe-standby action (zero net power), preserving
safety at the cost of one lost dispatch cycle.

\section{Observability Outputs}

Each DC3S step emits a structured certificate containing:
\begin{itemize}
  \item Proposed and safe actions (charge/discharge MW).
  \item Reliability score $w_t$ and its component features.
  \item CQR interval width and adaptive multiplier.
  \item Solver status and objective value.
  \item SHA-256 certificate hash for audit traceability.
\end{itemize}

These outputs are logged to \texttt{reports/hil/hil\_certificate\_audit.csv}
(Chapter~\ref{ch:hil}) and can be streamed to an observability backend
(e.g., Prometheus/Grafana) for real-time monitoring.

\section{Deployment Implications}

The sub-millisecond median latency has three practical consequences:

\begin{enumerate}
  \item \textbf{Frequency regulation}: DC3S can serve 4-second AGC
    signals with $>$99.99\% on-time delivery.
  \item \textbf{Fleet scaling}: At 0.042\,ms per battery, a single
    core can certify $\sim$700 batteries per 30-second cycle, enabling
    fleet-level deployment without parallelisation
    (cf.\ Chapter~\ref{ch:compositional-safety-fleets}).
  \item \textbf{Edge deployment}: The combined latency and memory
    footprint are compatible with ARM-based edge gateways, eliminating
    cloud round-trip dependencies.
\end{enumerate}

\section{Summary}

The DC3S safety pipeline completes in a median of 0.037\,ms per step,
with worst-case latency of 2.4\,ms.  The full pipeline including
forecast, OQE, RUI, SAF, and certificate generation fits within
$\sim$12\,MB RAM.  These figures confirm that the safety guarantee
is computationally free in the context of grid-scale dispatch cycles,
removing latency as a barrier to deployment.
